Cardiovascular diseases (CVDs) accounted for approximately 19.8 million deaths in 2022 \cite{who2025}, making their identification and prevention a major public health priority. Although coronary heart disease is traditionally diagnosed through clinical evaluation and diagnostic testing \cite{chestpain2023}, population-level health surveys provide an added perspective on the demographic, lifestyle, and health-related factors associated with the disease. Machine learning enables the exploration of these relationships, although the quality of the resulting classification depends considerably on the characteristics of the available data. Missing values, redundant features, and, in particular, class imbalance present important challenges, as underrepresented positive cases can be overlooked despite high overall accuracy \cite{brfss2026}. Against this background, the project investigates the classification of self-reported coronary heart disease using BRFSS data, with particular emphasis on data preprocessing, feature representation, and comparative evaluation of linear classification models.